% GENERATED FILE. Edit canonical lessons, then run npm run generate:books.
% canonical-source-hash: fnv1a64:c3637894a6e5e616
% canonical-lessons: UR-C51-bhulna, UR-C51-chahna, UR-C51-lana, UR-C51-bhejna, UR-C51-ginna

\chapter{To forget, To want, To bring, To send, To count}
\label{ch:ur-to-forget-to-want-to-bring-to-send-to-count}
\hlchaptermodality{51}

\begin{chapteropening}
\textbf{By the end of this chapter:} \emph{I can say to forget, to want, to bring, to send and to count.}

Complete the last lesson of chapter 51: I can say to forget, to want, to bring, to send and to count.
\end{chapteropening}

\section[\texorpdfstring{\ur{بھولنا} (bhūlnā)}{بھولنا (bhūlnā)}]{\ur{بھولنا} (bhūlnā) — to forget}

\label{lesson:UR-C51-bhulna}



\begin{quote}
Before the new one: say the Urdu for to live, then the Urdu for to learn.
\end{quote}

\subsection*{You'll want to know: \ur{بھولنا}}
\textbf{\ur{بھولنا}} — \emph{bhūlnā} — \textquotedblleft{}to forget\textquotedblright{}.

\begin{cousinweb}
\textbf{\ur{بھولنا}} (\emph{bhūlnā}) is a tadbhava, a word inherited and reshaped by sound change. It comes from \emph{*bhull}, \textquotedblleft{}err, forget\textquotedblright{}, a form CDIAL reconstructs from its descendants (CDIAL 9538).
\end{cousinweb}

\begin{grammarlens}[title={What you've built}]
One more word for this chapter.
\end{grammarlens}

\subsection*{Your turn}
\begin{itemize}
\raggedright
  \item \emph{Say it:} \emph{bhūlnā}
  \item \emph{Say it:} \emph{bhūlnā}, once more
  \item \emph{Say it:} say \emph{bhūlnā}
  \item \emph{Recall:} say the Urdu for to live, then the Urdu for to learn, then say \emph{bhūlnā} again
\end{itemize}

\begin{tcolorbox}[breakable,colback=teal!4,colframe=teal!35!black,title={Before you move on}]
What does \ur{بھولنا} mean? (\textquotedblleft{}To forget\textquotedblright{}.) Say it once more.
\end{tcolorbox}
\section[\texorpdfstring{\ur{چاہنا} (chāhnā)}{چاہنا (chāhnā)}]{\ur{چاہنا} (chāhnā) — to want}

\label{lesson:UR-C51-chahna}



\begin{quote}
Before the new one: say the Urdu for to learn, then the Urdu for to forget.
\end{quote}

\subsection*{You'll want to know: \ur{چاہنا}}
\textbf{\ur{چاہنا}} — \emph{chāhnā} — \textquotedblleft{}to want\textquotedblright{}.

\begin{cousinweb}
\textbf{\ur{چاہنا}} (\emph{chāhnā}) is a tadbhava, a word inherited and reshaped by sound change. It comes from \emph{*cāh}, \textquotedblleft{}see, look for, desire\textquotedblright{}, a form CDIAL reconstructs from its descendants (CDIAL 4775). CDIAL calls its deeper source uncertain.
\end{cousinweb}

\begin{grammarlens}[title={What you've built}]
One more word for this chapter.
\end{grammarlens}

\subsection*{Your turn}
\begin{itemize}
\raggedright
  \item \emph{Say it:} \emph{chāhnā}
  \item \emph{Say it:} \emph{chāhnā}, once more
  \item \emph{Say it:} say \emph{chāhnā}
  \item \emph{Recall:} say the Urdu for to learn, then the Urdu for to forget, then say \emph{chāhnā} again
\end{itemize}

\begin{tcolorbox}[breakable,colback=teal!4,colframe=teal!35!black,title={Before you move on}]
What does \ur{چاہنا} mean? (\textquotedblleft{}To want\textquotedblright{}.) Say it once more.
\end{tcolorbox}
\section[\texorpdfstring{\ur{لانا} (lānā)}{لانا (lānā)}]{\ur{لانا} (lānā) — to bring}

\label{lesson:UR-C51-lana}



\begin{quote}
Before the new one: say the Urdu for to forget, then the Urdu for to want.
\end{quote}

\subsection*{You'll want to know: \ur{لانا}}
\textbf{\ur{لانا}} — \emph{lānā} — \textquotedblleft{}to bring\textquotedblright{}.

\begin{cousinweb}
\textbf{\ur{لانا}} (\emph{lānā}) is inherited. CDIAL lists it under Sanskrit \emph{labhate}, \textquotedblleft{}catches, takes\textquotedblright{}, crossed with forms of \emph{dadāti}, \textquotedblleft{}gives\textquotedblright{} (CDIAL 10948): to bring is to take and give.
\end{cousinweb}

\begin{grammarlens}[title={What you've built}]
One more word for this chapter.
\end{grammarlens}

\subsection*{Your turn}
\begin{itemize}
\raggedright
  \item \emph{Say it:} \emph{lānā}
  \item \emph{Say it:} \emph{lānā}, once more
  \item \emph{Say it:} say \emph{lānā}
  \item \emph{Recall:} say the Urdu for to forget, then the Urdu for to want, then say \emph{lānā} again
\end{itemize}

\begin{tcolorbox}[breakable,colback=teal!4,colframe=teal!35!black,title={Before you move on}]
What does \ur{لانا} mean? (\textquotedblleft{}To bring\textquotedblright{}.) Say it once more.
\end{tcolorbox}
\section[\texorpdfstring{\ur{بھیجنا} (bhejnā)}{بھیجنا (bhejnā)}]{\ur{بھیجنا} (bhejnā) — to send}

\label{lesson:UR-C51-bhejna}



\begin{quote}
Before the new one: say the Urdu for to want, then the Urdu for to bring.
\end{quote}

\subsection*{You'll want to know: \ur{بھیجنا}}
\textbf{\ur{بھیجنا}} — \emph{bhejnā} — \textquotedblleft{}to send\textquotedblright{}.

\begin{cousinweb}
\textbf{\ur{بھیجنا}} (\emph{bhejnā}) is a tadbhava, a word inherited and reshaped by sound change. It comes from \emph{*bhejj}, \textquotedblleft{}send\textquotedblright{}, a form CDIAL reconstructs from its descendants (CDIAL 9603). CDIAL rejects a derivation from \emph{visarjayati}, \textquotedblleft{}sends off\textquotedblright{}, and offers two possible sources, neither certain.
\end{cousinweb}

\begin{grammarlens}[title={What you've built}]
One more word for this chapter.
\end{grammarlens}

\subsection*{Your turn}
\begin{itemize}
\raggedright
  \item \emph{Say it:} \emph{bhejnā}
  \item \emph{Say it:} \emph{bhejnā}, once more
  \item \emph{Say it:} say \emph{bhejnā}
  \item \emph{Recall:} say the Urdu for to want, then the Urdu for to bring, then say \emph{bhejnā} again
\end{itemize}

\begin{tcolorbox}[breakable,colback=teal!4,colframe=teal!35!black,title={Before you move on}]
What does \ur{بھیجنا} mean? (\textquotedblleft{}To send\textquotedblright{}.) Say it once more.
\end{tcolorbox}
\section[\texorpdfstring{\ur{گننا} (ginnā)}{گننا (ginnā)}]{\ur{گننا} (ginnā) — to count}

\label{lesson:UR-C51-ginna}



\begin{quote}
Before the new one: say the Urdu for to bring, then the Urdu for to send.
\end{quote}

\subsection*{You'll want to know: \ur{گننا}}
\textbf{\ur{گننا}} — \emph{ginnā} — \textquotedblleft{}to count\textquotedblright{}.

\begin{cousinweb}
\textbf{\ur{گننا}} (\emph{ginnā}) is a tadbhava, a word inherited from Sanskrit and reshaped by sound change. It comes from \emph{gaṇayati}, \textquotedblleft{}counts\textquotedblright{} (CDIAL 3993).
\end{cousinweb}

\begin{grammarlens}[title={What you've built}]
One more word for this chapter.
\end{grammarlens}

\subsection*{Your turn}
\begin{itemize}
\raggedright
  \item \emph{Say it:} \emph{ginnā}
  \item \emph{Say it:} \emph{ginnā}, once more
  \item \emph{Say it:} say \emph{ginnā}
  \item \emph{Recall:} say the Urdu for to bring, then the Urdu for to send, then say \emph{ginnā} again
\end{itemize}

\begin{tcolorbox}[breakable,colback=teal!4,colframe=teal!35!black,title={Before you move on}]
What does \ur{گننا} mean? (\textquotedblleft{}To count\textquotedblright{}.) Say it once more.
\end{tcolorbox}
